\section{与 SAM3 Agent 的对比}

\subsection{显式构建 vs. 自主涌现}

上一期介绍的 SAM3 Agent 也实现了 Visual Prompting，但方式截然不同：

\begin{table}[H]
\centering
\begin{tabular}{p{3cm}p{5cm}p{5cm}}
\toprule
\textbf{维度} & \textbf{SAM3 + VLM} & \textbf{Gemini Agentic Vision} \\
\midrule
视觉工具 & SAM3（专用分割模型） & Code Execution（通用） \\
架构 & 显式构建的 Agent Pipeline & Post-Training 激励出的行为 \\
视觉草图 & 由 SAM3 生成 Mask 叠加图 & 模型自主写代码生成 \\
BBox 来源 & 外部模型/工具 & 模型原生能力 \\
灵活性 & 受限于定义的 Tool & 理论上可完成任意图像操作 \\
\bottomrule
\end{tabular}
\caption{SAM3 Agent 与 Gemini Agentic Vision 对比}
\end{table}

\subsection{Visual Prompting 的演进}

两种方案都产生了"视觉草图"（Visual Scratchpad），但来源不同：

\begin{itemize}
    \item SAM3 Agent：人类\textbf{显式设计}了分割→标注→返回 VLM 的流程
    \item Agentic Vision：模型\textbf{自主决定}需要什么样的视觉草图来辅助推理
\end{itemize}

\subsection{本章小结}

从显式构建到自主涌现，Agentic Vision 代表了 Visual Prompting 技术的进化方向——让模型自己决定如何"看"图，而非人类预设好观察方式。


